\section{Introduction}
\label{sec:intro}

Class-incremental learning (CIL) remains a fundamental challenge in continual visual learning. A model must sequentially acquire new visual categories without access to previous training data, and ultimately classify test samples into all seen classes without task identities~\cite{rebuffi2017icarl,van2019three}. This creates a persistent stability-plasticity dilemma: the model must simultaneously develop task-specialized representations for fine-grained discrimination (plasticity) while preserving cross-task knowledge to prevent forgetting (stability). The tension is particularly acute at test time, where the model must navigate a continually expanding label space with no explicit task signal.

With large-scale pretrained Vision Transformers (ViTs), structural expansion, allocating per-task modules on a frozen backbone, has become a predominant paradigm in CIL. By isolating task-specific parameters, expansion methods inherently avoid interference with previously learned knowledge. Nevertheless, this consensus masks a practical bottleneck. Current expansion strategies incur escalating \emph{parameter growth} (e.g., prompt pools, adapter chains, or full experts growing linearly with the number of tasks) and escalating \emph{inference cost} (e.g., requiring traversal or selection across all stored modules at test time). As the task horizon lengthens, both costs become increasingly burdensome.

Beyond these cost concerns, existing structural methods share a more fundamental weakness. Their test-time inference depends on an external routing decision whose failure is catastrophic and irrecoverable. Prompt- and adapter-based methods~\cite{wang2022learning,ease2024,wang2025integrating} select task-specific modules through key-query matching or hard module selection. A single misrouting redirects the prediction to an irrelevant label space. Since logits from different modules are mutually uncalibrated, no post-hoc fusion can meaningfully correct the error. Prototype-based methods~\cite{ranpac2023,fecam2023} circumvent routing entirely by operating in a single globally consistent feature space, yet they sacrifice the fine-grained discrimination of task-adapted experts. A structural dilemma thus remains unresolved: \emph{specialized experts require brittle hard routing}, while \emph{routing-free classifiers lack expert-level discrimination}.

SPIE (Scalable Posterior Inference with Low-Rank Experts) resolves this trade-off by uniting specialist-level discrimination with routing-free inference. The key idea is threefold: construct per-task compact experts for local discriminability, maintain a globally calibrated prototype space for task routing, and fuse their predictions through a soft posterior mechanism. This mechanism eliminates the brittleness of hard task selection while preserving expert-level plasticity.

\textit{First}, a lightweight stable-plastic decomposition couples a globally calibrated prototype space (maintained on a frozen pretrained ViT backbone) with per-task compact experts. Each expert learns merely a set of scaling vectors that modulate a frozen low-rank basis derived from the pretrained weights. This yields less than \textbf{0.1\%} parameter growth per task while retaining strong local discriminability.

\textit{Second}, a low-rank distributional prototype replay mechanism stores class-wise feature statistics in a compact diagonal-plus-low-rank form. It replays them as synthetic features to keep the globally calibrated prototype space aligned across all tasks, without storing raw data.

\textit{Third}, a soft posterior inference procedure uses the calibrated prototype vectors as a unified prototype bank to estimate a task posterior distribution $q(t \mid x)$. It then fuses expert-specific class probabilities weighted by this posterior in probability space. This design eliminates hard task selection while preventing direct comparison of uncalibrated logits.

Our contributions are three-fold:
\begin{enumerate}[leftmargin=*,nosep]
    \item We propose a stable-plastic CIL framework that couples a globally calibrated prototype space with compact per-task experts and low-rank distributional prototype storage. This constrains parameter growth to below 0.1\% of the backbone per incremental task.
    \item We introduce a soft posterior inference mechanism that derives task-level routing signals from a unified prototype bank and fuses expert predictions in probability space. This eliminates the brittle hard task-ID dependence that afflicts existing structural methods.
    \item We conduct extensive experiments across multiple standard CIL benchmarks. Results demonstrate that SPIE achieves strong performance with substantially lower parameter and inference costs than prior structural expansion methods.
\end{enumerate}
